\textbf{This section formalizes DWTS as a replayable inverse problem and specifies the modeling objects that close our pipeline.}
We define the weekly active set, the observable judges' score signal, and the realized elimination outcomes, then introduce the latent fan vote share $p_{i,t}$ as the primary unknown.
Our inference module outputs point estimates $\hat p_{i,t}$ together with uncertainty summaries, which feed a unified rule replay engine; this engine generates counterfactual elimination paths used for rule comparison, case-based attribution, and the validation of a proposed new voting system.

% ===========================问题分解与分析=============================
\subsection{Problem Decomposition}\label{sec:problem_decomposition}

\textbf{We decompose the prompt into four linked tasks with shared inputs and a common replay-based evaluation layer.}
\begin{itemize}
    \item \textbf{Task A: Fan vote inference with uncertainty.}
    \textbf{Input:} weekly judges' scores and the realized eliminated set;
    \textbf{Output:} weekly fan vote shares $p_{i,t}$ (and uncertainty bands) over the active set;
    \textbf{Validation:} elimination-consistency under the assumed rule, plus identifiability diagnostics (feasible-range widths and margin-based sensitivity).

    \item \textbf{Task B: Voting-rule comparison via replay.}
    \textbf{Input:} judges' scores and inferred $p_{i,t}$;
    \textbf{Output:} counterfactual week-by-week eliminations under percent-based vs.\ rank-based aggregation (and variants);
    \textbf{Validation:} season-level and week-level outcome divergence measures and ``fan bias'' indicators.

    \item \textbf{Task C: Controversial contestants and judges' save attribution.}
    \textbf{Input:} inferred $p_{i,t}$ and replay engine;
    \textbf{Output:} case-study narratives supported by counterfactual trajectories with/without bottom-two judges' save;
    \textbf{Validation:} whether the mechanism explains observed controversies and how sensitive conclusions are to vote uncertainty.

    \item \textbf{Task D: Mechanism analysis and new system design.}
    \textbf{Input:} panel features (celebrity attributes, professional dancer identifiers) and inferred votes;
    \textbf{Output:} mixed-effects estimates for drivers of judges' scoring versus fan preference, and a parameterized voting rule;
    \textbf{Validation:} historical replay, robustness checks, and sensitivity analysis over hyperparameters and rule-switch assumptions.
\end{itemize}


% ==========================关键假设=============================
\subsection{Key Assumptions}\label{subsec:assumptions}

We adopt a small, testable assumption budget to make the inverse problem well-posed while keeping conclusions falsifiable.
\begin{enumerate}
    \item \textbf{Share normalization.} Fan votes are modeled as within-week shares over the active set.
    \item \textbf{Rule-governed elimination.} Eliminations follow the documented aggregation rule (percent-based or rank-based), optionally augmented by bottom-two judges' save where applicable.
    \item \textbf{Active-set definition and special episodes.} Post-exit ``zeros'' represent inactivity and are excluded from the active set; special episodes are explicitly flagged and may be accommodated via slack rather than forced exact compliance. When a counterfactual replay keeps an observed-exit contestant active, the required post-exit inputs are filled using a conservative, minimal-information imputation; this is used only to extend counterfactual trajectories and is not intended as a performance forecast.
    \item \textbf{Judges as observed signal.} Judges' scores provide an observed performance signal taken as exogenous to the voting mechanism.
    \item \textbf{Temporal smoothness.} Fan preference evolves smoothly week-to-week, motivating temporal regularization.
\end{enumerate}
The season at which the show returns to rank-based aggregation is treated as an explicit assumption and is stress-tested using adjacent cutoffs.

% =========================记号=============================
\subsection{Key Notations}\label{subsec:notations}
Table~\ref{tab:notations} summarizes indices, sets, observed quantities, latent fan vote shares, and rule-dependent combined scores used throughout the inference and replay pipeline.
\begin{table}[H]
\centering
\small
\begin{tabular}{ll}
\toprule
Symbol & Definition \\
\midrule
$s$ & Season index, $s=1,\dots,34$ \\
$t$ & Week index within season $s$ \\
$i$ & Contestant (celebrity--pro pair) index \\
$A_{s,t}$ & Active contestants in season $s$, week $t$ \\
$E_{s,t}$ & Eliminated set at end of week $t$; $m_{s,t}=|E_{s,t}|$ \\
$j_{i,t,k}$ & Score given by judge $k$ to contestant $i$ in week $t$ \\
$J_{i,t}$ & Total judges' score: $J_{i,t}=\sum_{k} j_{i,t,k}$ (skip missing) \\
$q^{J}_{i,t}$ & Judges' share: $q^{J}_{i,t}=J_{i,t}\big/\sum_{r\in A_{s,t}}J_{r,t}$ \\
$p_{i,t}$ & Fan vote share (latent), $\sum_{i\in A_{s,t}}p_{i,t}=1$ \\
$\alpha$ & Weight on judges under percent rule (default $\alpha=0.5$; stress-tested) \\
$S^{(P)}_{i,t}$ & Percent combined score: $S^{(P)}_{i,t}=\alpha q^{J}_{i,t}+(1-\alpha)p_{i,t}$ \\
$R^{J}_{i,t},R^{V}_{i,t}$ & Judges rank and fan rank (descending) \\
$R^{(R)}_{i,t}$ & Rank combined score: $R^{(R)}_{i,t}=R^{J}_{i,t}+R^{V}_{i,t}$ \\
$\xi_t$ & Slack for elimination-consistency; penalized in vote inference \\
$\lambda$ & Penalty weight on slacks in vote inference objective \\
\bottomrule
\end{tabular}
\caption{Core notations used throughout the paper.}
\label{tab:notations}
\end{table}

\subsection{Tasks and Evaluation KPIs}\label{subsec:tasks_and_kpis}

Our workflow consists of four tightly coupled tasks. \textbf{Task A} reconstructs latent weekly fan vote shares $p_{i,t}$
for each active contestant by solving a season-level constrained optimization problem that enforces elimination consistency
(with slack) under the documented mechanics\cite{bubeck2015convex,boyd2004convex}. \textbf{Tasks B/C} use these inferred votes to replay each season under
alternative aggregation rules (rank-based vs.\ percent-based, optionally with a bottom-two judges' save), producing
counterfactual eliminations and placements. \textbf{Task D1} fits mixed-effects models to contrast the drivers of judges'
evaluations versus audience support (celebrity attributes and professional partners). \textbf{Task D2} proposes an
implementable voting rule and selects parameters via replay-based stress tests.

To evaluate models and rule variants consistently, we use the following KPIs and slices.

\textbf{KPI1: Elimination Set Agreement (Jaccard Similarity).}
On each season-week $(s,t)$, we compute the replayed eliminated set $\hat{E}_{s,t}$ under a specified rule
and compare it to the observed eliminated-at-end-of-week set $E_{s,t}$ parsed from results. We measure
set-level agreement by the Jaccard similarity
\[
J_{s,t}=\frac{|E_{s,t}\cap \hat{E}_{s,t}|}{|E_{s,t}\cup \hat{E}_{s,t}|}.
\]
KPI1 is reported as the week-level mean of $J_{s,t}$ over all counted season-week instances.
Weeks with no elimination yield $E_{s,t}=\varnothing$ (and possibly $\hat{E}_{s,t}=\varnothing$); we adopt the
convention $J_{s,t}=1$ when both sets are empty and include such weeks in the mean.

\paragraph{KPI3: Rule Divergence Rate (Rank vs.\ Percent).}
Holding the same vote input (e.g., BL-0 proxy or inferred votes), KPI3 measures how often two aggregation schemes disagree
on eliminations:
\begin{equation}
\text{KPI3}=\frac{\sum_{(s,t)\in\Omega}\mathbf{1}\{\widehat{E}^{\text{rank}}_{s,t}\neq \widehat{E}^{\text{percent}}_{s,t}\}}{|\Omega|},
\end{equation}
where $\Omega$ denotes eligible $(s,t)$. KPI3 is reported descriptively as ``impact magnitude,'' not as an optimization objective.

\paragraph{KPI2: Vote Identifiability (interval tightness).}
Because votes are unobserved, vote inference is not always identifiable.\cite{benning2018regularization,kaipio2005statistical} For each $(s,t,i)$, we compute a feasible interval
$[p^{\min}_{i,t},p^{\max}_{i,t}]$ over fan vote shares consistent with elimination constraints (and with a near-optimal shell
around the Task~A objective). We define an identifiability score as
\begin{equation}
ID_{i,t}=\frac{1}{p^{\max}_{i,t}-p^{\min}_{i,t}+\varepsilon},
\end{equation}
and summarize by quantiles across contestants, weeks, and seasons. KPI2 flags weeks and contestants where multiple vote
allocations explain the observed elimination equally well.

\paragraph{Reproducible baseline (BL-0).}
Before introducing inferred votes, we report KPI1/KPI3 under BL-0 as an audit baseline. In BL-0, the rank scheme reduces to
judges' ranks only (uniform fan ranks add a constant), and the percent scheme uses judges' percent plus a uniform fan share
$1/n_{\text{active}}$. BL-0 validates denominator handling and replay mechanics, but it is not used to claim audience preference estimates.
